\documentclass[../../main.tex]{subfiles}% 注意这里的文件路径不能用 ./main.tex ,否则用latexmk编译子文件会报错
\graphicspath{{\subfix{./image/}}{\subfix{../../image/}}} % 指定图片引用路径,后续可以直接使用图片文件名
% 如果图片存放在上级目录,则使用 ../../image/ 作为备用路径

% 例如:
% \begin{figure}[H]
% \centering
% \includegraphics[width=0.3\textwidth]{图.png}
% \caption{}
% \label{figure:图}
% \end{figure}
% 注意:上述\label{}一定要放在\caption{}之后,否则引用图片序号会只会显示??.

\begin{document}

\section{测度}\label{section:ch01-测度}

\subsection{一般性质}\label{subsection:ch01-测度一般性质}

\begin{definition}\label{definition:ch01-集合函数性质}
设 $X$ 为一集合，$\mathcal{F}\subseteq\mathcal{P}(X)$ 为一集族，满足 $\varnothing \in\mathcal{F}$。映射 $\varphi:\mathcal{F}\to[0,\infty]$ 称为 $\mathcal{F}$ 上的集合函数。它称为：
\begin{enumerate}[(i)]
\item 单调的，若对任意 $E,F\in\mathcal{F}$，$E\subseteq F$ 蕴含 $\varphi(E)\leqslant\varphi(F)$；
\item 可加的，若对任意 $n\in\mathbb{N}$ 与任意有限不交族 $\{E_1,\cdots ,E_n\}\subseteq\mathcal{F}$，当 $\bigcup_{k=1}^nE_k\in\mathcal{F}$ 时有
\begin{align*}
\varphi\left(\bigcup_{k=1}^nE_k\right)=\sum_{k=1}^n\varphi(E_k);
\end{align*}
\item $\sigma$--次可加的，若对任意序列 $\{E_k\}\subseteq\mathcal{F}$，当 $\bigcup_{k=1}^\infty E_k\in\mathcal{F}$ 时有
\begin{align*}
\varphi\left(\bigcup_{k=1}^\infty E_k\right)\leqslant\sum_{k=1}^\infty\varphi(E_k);
\end{align*}
\item $\sigma$--可加的，若对任意不交序列 $\{E_k\}\subseteq\mathcal{F}$，当 $\bigcup_{k=1}^\infty E_k\in\mathcal{F}$ 时有
\begin{align*}
\varphi\left(\bigcup_{k=1}^\infty E_k\right)=\sum_{k=1}^\infty\varphi(E_k).
\end{align*}
\end{enumerate}
\end{definition}

% SUPPLEMENT: parent=definition:ch01-集合函数性质; kind=motivation; 说明四种性质的动机、$\varnothing\in\mathcal{F}$ 与"并仍在 $\mathcal{F}$ 中"两条假设的作用，以及四者之间的强弱关系
\begin{remark}\phantomsection\label{remark:ch01-集合函数性质}
\begin{enumerate}[(i)]
\item 引入动机。测度的构造分两步：先在便于计算体积、但不便作可数运算的集族上定义集合函数，再把它延拓到 $\sigma$--代数（参见\cref{theorem:ch01-测度延拓}）。第一步能检验的只是有限可加性、单调性等弱性质，因此需要把这些性质作为独立概念区分开；而 $\sigma$--可加性正是测度的定义性要求（\cref{definition:ch01-测度}）。
\item 假设的作用。$\varnothing \in\mathcal{F}$ 保证空集上的值有意义，从而 $\sigma$--可加性可以排除 $\varphi(\varnothing )=\infty$ 这一退化情形；条件 $\bigcup_{k=1}^nE_k\in\mathcal{F}$（相应地 $\bigcup_{k=1}^\infty E_k\in\mathcal{F}$）保证等式左端有定义，因为 $\varphi$ 只在 $\mathcal{F}$ 上取值。半代数（\cref{definition:ch01-半代数代数}）一般不对并封闭，故该条件不能省去。
\item 强弱关系。在 $\mathcal{F}$ 对有限并、差集封闭（例如 $\mathcal{F}$ 为 $\sigma$--代数）的前提下，$\sigma$--可加性蕴含可加性，也蕴含 $\sigma$--次可加性，二者又共同给出单调性（见下方的证明）。反之，$\sigma$--次可加性不蕴含 $\sigma$--可加性：$[0,1]$ 上的 Lebesgue 外测度 $\lambda^*$ 是 $\sigma$--次可加的，但它不是 $\mathcal{P}([0,1])$ 上的测度。
\item 边界情形。若 $\mathcal{F}=\{\varnothing \}$，则四条性质都只涉及空集，其中 $\sigma$--可加性给出 $\varphi(\varnothing )=\sum_{k=1}^\infty\varphi(\varnothing )$，故 $\varphi(\varnothing )\in\{0,\infty\}$；可见即使论域退化为单点族，$\varphi(\varnothing )=0$ 仍是一条独立要求。
\end{enumerate}
\end{remark}

% SUPPLEMENT: parent=definition:ch01-集合函数性质; kind=proof; 补出第 (iii) 项所列的蕴含关系：$\sigma$--可加性 $\Rightarrow$ 可加性、$\sigma$--可加性 $\Rightarrow$ $\sigma$--次可加性、以及 $\sigma$--可加性 $\Rightarrow$ 单调性
\begin{proof}
设 $\mathcal{F}$ 对有限并与差集封闭，$\varnothing \in\mathcal{F}$，且 $\varphi:\mathcal{F}\to[0,\infty]$ 为 $\sigma$--可加的。

可加性：设 $\{E_1,\cdots ,E_n\}\subseteq\mathcal{F}$ 不交且 $E\triangleq\bigcup_{k=1}^nE_k\in\mathcal{F}$。令 $E_k\triangleq\varnothing$（$k\geqslant n+1$）。则 $\{E_k\}_{k\in\mathbb{N}}$ 是 $\mathcal{F}$ 中的不交序列且 $\bigcup_{k=1}^\infty E_k=E\in\mathcal{F}$，故由 $\sigma$--可加性
\begin{align*}
\varphi(E)=\sum_{k=1}^\infty\varphi(E_k)=\sum_{k=1}^n\varphi(E_k),
\end{align*}
其中末一步用到 $\varphi(\varnothing )=0$。

$\sigma$--次可加性：设 $\{E_k\}\subseteq\mathcal{F}$ 且 $E\triangleq\bigcup_{k=1}^\infty E_k\in\mathcal{F}$。令
\begin{align*}
F_1\triangleq E_1,\qquad F_k\triangleq E_k\setminus\bigcup_{j=1}^{k-1}E_j\quad(k\geqslant2).
\end{align*}
由 $\mathcal{F}$ 对有限并与差集封闭得 $\{F_k\}\subseteq\mathcal{F}$；序列 $\{F_k\}$ 不交，$F_k\subseteq E_k$，且 $\bigcup_{k=1}^\infty F_k=E$。于是由 $\sigma$--可加性与 $\varphi\geqslant0$
\begin{align*}
\varphi(E)=\sum_{k=1}^\infty\varphi(F_k)\leqslant\sum_{k=1}^\infty\varphi(E_k).
\end{align*}

单调性：设 $E,F\in\mathcal{F}$，$E\subseteq F$。则 $F\setminus E\in\mathcal{F}$，$F=E\cup(F\setminus E)$ 为不交并，故由可加性与 $\varphi\geqslant0$ 得 $\varphi(F)=\varphi(E)+\varphi(F\setminus E)\geqslant\varphi(E)$。
\end{proof}

\begin{definition}\label{definition:ch01-测度}
设 $(X,\mathcal{A})$ 为一可测空间。集合函数 $\mu:\mathcal{A}\to[0,\infty]$ 称为 $\mathcal{A}$ 上（或 $X$ 上）的一个（正）测度，若 $\mu(\varnothing )=0$ 且 $\mu$ 为 $\sigma$--可加的。三元组 $(X,\mathcal{A},\mu)$ 称为一个测度空间。

测度 $\mu$ 称为有限的，若 $\mu(X)<\infty$；它称为 $\sigma$--有限的，若存在序列 $\{E_k\}\subseteq\mathcal{A}$ 使得 $X=\bigcup_{k=1}^\infty E_k$ 且对所有 $k\in\mathbb{N}$ 有 $\mu(E_k)<\infty$。测度空间 $(X,\mathcal{A},\mu)$ 称为有限的（相应地，$\sigma$--有限的），若 $\mu$ 为有限的（相应地，$\sigma$--有限的）。

若 $\mu(X)=1$，则测度 $\mu$ 称为概率测度，空间 $(X,\mathcal{A},\mu)$ 称为概率空间。$X$ 上概率测度的全体记为 $\mathcal{P}(X)$。

设 $\mathcal{F}\subseteq\mathcal{P}(X)$ 为任意集族。映射 $\mu:\mathcal{F}\to[0,\infty]$ 称为 $\mathcal{F}$ 上的一个测度，若当 $\varnothing \in\mathcal{F}$ 时 $\mu(\varnothing )=0$，且对任意不交序列 $\{E_k\}\subseteq\mathcal{F}$，当 $\bigcup_{k=1}^\infty E_k\in\mathcal{F}$ 时有
\begin{align*}
\mu\left(\bigcup_{k=1}^\infty E_k\right)=\sum_{k=1}^\infty\mu(E_k).
\end{align*}
$\mathcal{F}$ 上的测度 $\mu$ 称为 $\sigma$--有限的，若存在序列 $\{E_k\}\subseteq\mathcal{F}$ 使得 $X=\bigcup_{k=1}^\infty E_k$ 且对每个 $k\in\mathbb{N}$ 有 $\mu(E_k)<\infty$。
\end{definition}

% SUPPLEMENT: parent=definition:ch01-测度; kind=motivation; 说明 $\mu(\varnothing)=0$ 与 $\sigma$--可加性各自的作用、取值于 $[0,\infty]$ 的必要性、$\sigma$--有限性的用途，以及把定义推广到任意集族 $\mathcal{F}$ 的理由
\begin{remark}\phantomsection\label{remark:ch01-测度动机}
\begin{enumerate}[(i)]
\item 引入动机。测度是"体积"概念的抽象：$\sigma$--可加性使测度对可数不交并保持求和，因而是积分理论（单调收敛、Fubini--Tonelli）得以运转的最低要求；$\mathcal{A}$ 为 $\sigma$--代数则保证定义域对定义中出现的可数运算封闭。
\item 假设的作用。$\mu(\varnothing )=0$ 不能由 $\sigma$--可加性推出：若 $\mu(\varnothing )=\infty$，则 $\mu(\varnothing )=\sum_{k=1}^\infty\mu(\varnothing )$ 仍成立，故必须单独要求。$\sigma$--可加性只对不交序列陈述，使用时总要把序列不交化（见\cref{proposition:ch01-测度性质}的证明）。
\item 取值于 $[0,\infty]$。允许 $\mu$ 取 $\infty$ 是不可省的：$X$ 上的计数测度与 $\mathbb{R}^N$ 上的 Lebesgue 测度一般不是有限测度（\cref{example:ch01-计数与Dirac}）。
\item $\sigma$--有限性的作用。它把无穷测度化归为可数个有限块上的论证：例如在\cref{proposition:ch01-可数基Borel正则}的证明中，先在有限测度 $\mu|_{A_n}$ 上得到外正则性与内正则性，再对 $n$ 求和。若只假设有限性，则 $\mathbb{R}^N$ 上的 Lebesgue 测度已被排除。
\item 推广到任意集族。定义的最后一段把测度概念移到不对并、补封闭的族 $\mathcal{F}$ 上，这是测度延拓定理的出发点：只需在 $\mathcal{F}$ 上验证不交可加性即可。此处 $\sigma$--有限性指用 $\mathcal{F}$ 中的集合覆盖 $X$，故它比 $\mathcal{A}$ 上的 $\sigma$--有限性更强。
\item 记号。定义中的 $\mathcal{P}(X)$ 在第三段表示概率测度族，在第一段（以及\cref{definition:ch01-集合函数性质}）表示幂集，两者按上下文区分。
\end{enumerate}
\end{remark}

\begin{remark}\phantomsection\label{remark:ch01-测度子空间}
设 $(X,\mathcal{A},\mu)$ 为一测度空间，$F\in\mathcal{A}$。$\mu$ 在迹 $\sigma$--代数 $\mathcal{A}\cap F$ 上的限制 $\mu|_{\mathcal{A}\cap F}$ 是由 $\mu$ 在 $\mathcal{A}\cap F$ 上诱导的测度。测度空间 $(F,\mathcal{A}\cap F,\mu|_{\mathcal{A}\cap F})$ 称为 $(X,\mathcal{A},\mu)$ 的一个测度子空间。
\end{remark}

% SUPPLEMENT: parent=remark:ch01-测度子空间; kind=proof; 原书只断言限制为测度，此处验证迹 $\sigma$--代数是 $F$ 上的 $\sigma$--代数、限制映射良定义且 $\sigma$--可加
\begin{proof}
记 $\mathcal{A}\cap F\triangleq\{E\cap F\mid E\in\mathcal{A}\}$。先证它是 $F$ 上的 $\sigma$--代数。由 $\varnothing =\varnothing \cap F$ 得 $\varnothing \in\mathcal{A}\cap F$。若 $A=E\cap F$，$E\in\mathcal{A}$，则 $A$ 在 $F$ 中的补集为 $F\setminus A=E^c\cap F\in\mathcal{A}\cap F$，因为 $E^c\in\mathcal{A}$。若 $\{A_k\}\subseteq\mathcal{A}\cap F$，$A_k=E_k\cap F$，则 $\bigcup_{k=1}^\infty A_k=\left(\bigcup_{k=1}^\infty E_k\right)\cap F\in\mathcal{A}\cap F$。故 $\mathcal{A}\cap F$ 是 $F$ 上的 $\sigma$--代数。

再证限制映射良定义。若 $E_1,E_2\in\mathcal{A}$ 满足 $E_1\cap F=E_2\cap F$，则 $\mu(E_1\cap F)=\mu(E_2\cap F)$，故 $A\mapsto\mu(A)$（$A\in\mathcal{A}\cap F$）不依赖于表示 $A=E\cap F$ 的选取。最后，设 $\{A_k\}\subseteq\mathcal{A}\cap F$ 不交，$A_k=E_k\cap F$，则 $\{A_k\}$ 在 $\mathcal{A}$ 中不交且 $\bigcup_{k=1}^\infty A_k\in\mathcal{A}\cap F$，故由 $\mu$ 的 $\sigma$--可加性
\begin{align*}
\mu|_{\mathcal{A}\cap F}\left(\bigcup_{k=1}^\infty A_k\right)=\mu\left(\bigcup_{k=1}^\infty A_k\right)=\sum_{k=1}^\infty\mu(A_k)=\sum_{k=1}^\infty\mu|_{\mathcal{A}\cap F}(A_k).
\end{align*}
又 $\mu|_{\mathcal{A}\cap F}(\varnothing )=\mu(\varnothing )=0$，故 $\mu|_{\mathcal{A}\cap F}$ 是 $F$ 上的测度。
\end{proof}

以下测度的性质可由定义直接推出。

\begin{proposition}\label{proposition:ch01-测度性质}
设 $(X,\mathcal{A},\mu)$ 为一测度空间。则：
\begin{enumerate}[(i)]
\item 对任意 $n\in\mathbb{N}$ 与任意有限不交族 $\{E_1,\cdots ,E_n\}\subseteq\mathcal{A}$，有
\begin{align*}
\mu\left(\bigcup_{k=1}^nE_k\right)=\sum_{k=1}^n\mu(E_k);
\end{align*}
\item 对任意 $E,F\in\mathcal{A}$，$E\subseteq F$ 蕴含 $\mu(E)\leqslant\mu(F)$；
\item 对任意序列 $\{E_k\}\subseteq\mathcal{A}$，有
\begin{align*}
\mu\left(\bigcup_{k=1}^\infty E_k\right)\leqslant\sum_{k=1}^\infty\mu(E_k);
\end{align*}
\item 对任意非减序列 $\{E_k\}\subseteq\mathcal{A}$，有
\begin{align*}
\mu\left(\bigcup_{k=1}^\infty E_k\right)=\lim_{k\to\infty}\mu(E_k);
\end{align*}
\item 对任意非增序列 $\{E_k\}\subseteq\mathcal{A}$ 且 $\mu(E_1)<\infty$，有
\begin{align*}
\mu\left(\bigcap_{k=1}^\infty E_k\right)=\lim_{k\to\infty}\mu(E_k).
\end{align*}
\end{enumerate}
\end{proposition}

% SUPPLEMENT: parent=proposition:ch01-测度性质; kind=proof; 原书称"可由定义直接推出"而未给证明，此处逐条补出，并指出第 (v) 款中有限性假设的使用位置
\begin{proof}
(i) 设 $\{E_1,\cdots ,E_n\}\subseteq\mathcal{A}$ 不交。令 $E_k\triangleq\varnothing$（$k\geqslant n+1$）。则 $\{E_k\}_{k\in\mathbb{N}}$ 为 $\mathcal{A}$ 中的不交序列，$\bigcup_{k=1}^\infty E_k=\bigcup_{k=1}^nE_k$，故由 $\mu$ 的 $\sigma$--可加性与 $\mu(\varnothing )=0$ 得
\begin{align*}
\mu\left(\bigcup_{k=1}^nE_k\right)=\sum_{k=1}^\infty\mu(E_k)=\sum_{k=1}^n\mu(E_k).
\end{align*}

(ii) 设 $E,F\in\mathcal{A}$，$E\subseteq F$。则 $F=E\cup(F\setminus E)$ 为不交并，且 $F\setminus E\in\mathcal{A}$，故由 (i) 与 $\mu\geqslant0$ 得 $\mu(F)=\mu(E)+\mu(F\setminus E)\geqslant\mu(E)$。

(iii) 设 $\{E_k\}\subseteq\mathcal{A}$，令 $F_1\triangleq E_1$、$F_k\triangleq E_k\setminus\bigcup_{j=1}^{k-1}E_j$（$k\geqslant2$）。则 $\{F_k\}\subseteq\mathcal{A}$ 不交，$F_k\subseteq E_k$，且 $\bigcup_{k=1}^\infty F_k=\bigcup_{k=1}^\infty E_k$。于是由 $\sigma$--可加性与 (ii)
\begin{align*}
\mu\left(\bigcup_{k=1}^\infty E_k\right)=\sum_{k=1}^\infty\mu(F_k)\leqslant\sum_{k=1}^\infty\mu(E_k).
\end{align*}

(iv) 设 $\{E_k\}\subseteq\mathcal{A}$ 非减。令 $E_0\triangleq\varnothing$、$F_k\triangleq E_k\setminus E_{k-1}$（$k\in\mathbb{N}$）。则 $\{F_k\}\subseteq\mathcal{A}$ 不交，$\bigcup_{k=1}^nF_k=E_n$，且 $\bigcup_{k=1}^\infty F_k=\bigcup_{k=1}^\infty E_k$。由 $\sigma$--可加性与部分和的定义，
\begin{align*}
\mu\left(\bigcup_{k=1}^\infty E_k\right)=\sum_{k=1}^\infty\mu(F_k)=\lim_{n\to\infty}\sum_{k=1}^n\mu(F_k)=\lim_{n\to\infty}\mu(E_n),
\end{align*}
其中末一步用 (i)。

(v) 设 $\{E_k\}\subseteq\mathcal{A}$ 非增且 $\mu(E_1)<\infty$。令 $G_k\triangleq E_1\setminus E_k$（$k\in\mathbb{N}$）。则 $\{G_k\}\subseteq\mathcal{A}$ 非减且 $\bigcup_{k=1}^\infty G_k=E_1\setminus\bigcap_{k=1}^\infty E_k$。由于 $E_k\subseteq E_1$、$\mu(E_1)<\infty$，由 (i) 得
\begin{align*}
\mu(G_k)=\mu(E_1)-\mu(E_k),\qquad\mu\left(E_1\setminus\bigcap_{j=1}^\infty E_j\right)=\mu(E_1)-\mu\left(\bigcap_{j=1}^\infty E_j\right),
\end{align*}
其中两个减式都有意义，因为被减项都不超过 $\mu(E_1)<\infty$。于是由 (iv) 与上式
\begin{align*}
\mu(E_1)-\mu\left(\bigcap_{k=1}^\infty E_k\right)=\lim_{k\to\infty}\mu(G_k)=\mu(E_1)-\lim_{k\to\infty}\mu(E_k),
\end{align*}
再减去有限数 $\mu(E_1)$ 即得 $\mu\left(\bigcap_{k=1}^\infty E_k\right)=\lim_{k\to\infty}\mu(E_k)$。
\end{proof}

% SUPPLEMENT: parent=proposition:ch01-测度性质; kind=motivation; 说明五款性质各自的用途，并给出第 (iii)、(v) 款中假设不可省的反例
\begin{remark}\phantomsection\label{remark:ch01-测度性质注}
\begin{enumerate}[(i)]
\item 引入动机。命题把 $\sigma$--可加性拆成几条便于使用的形式：(i) 是有限可加性，(ii) 是单调性，(iii) 是可数次可加性，(iv) 是自下而上的连续性（下连续），(v) 是自上而下的连续性（上连续）。它们构成测度论中一切极限论证的基本工具。
\item 假设的作用。(v) 中 $\mu(E_1)<\infty$ 用于作减法 $\mu(E_1)-\mu(E_k)$：非增序列的测度可能恒为 $\infty$，此时 $\lim_k\mu(E_k)=\infty$ 而 $\mu\left(\bigcap_kE_k\right)$ 可以有限。取 $X=\mathbb{N}$、$\mu$ 为计数测度、$E_k=\{n\in\mathbb{N}\mid n\geqslant k\}$，则 $\mu(E_k)=\infty$、$\bigcap_{k=1}^\infty E_k=\varnothing$，故 $\mu\left(\bigcap_kE_k\right)=0\neq\infty=\lim_k\mu(E_k)$。(iii) 中的不等号可取严格：取 $X=[0,1]$、$\mu=\lambda$、$E_1=E_2=[0,1]$，则 $\mu(E_1\cup E_2)=1<2=\mu(E_1)+\mu(E_2)$。
\item 边界情形。(iv) 与 (v) 中的极限都在 $[0,\infty]$ 中取，故允许取 $\infty$；(ii) 说明测度是单调集函数，从而\cref{definition:ch01-集合函数性质}意义下的单调性对测度自动成立。可加性与连续性并不互相蕴含：外测度是 $\sigma$--次可加且单调的，但一般不可加。
\end{enumerate}
\end{remark}

\begin{example}\label{example:ch01-计数与Dirac}
(i) 设 $X$ 为一集合。映射 $\mu^\#:\mathcal{P}(X)\to[0,\infty]$ 定义为
\begin{align*}
\mu^\#(E)\triangleq
\begin{cases}
E\text{ 的元素个数}, & \text{若 }E\text{ 为有限集},\\
\infty, & \text{否则},
\end{cases}
\end{align*}
是一个测度，称为计数测度。

(ii) 设 $X\neq\varnothing$，$x\in X$。映射 $\delta_x:\mathcal{P}(X)\to[0,\infty)$ 定义为
\begin{align*}
\delta_x(E)\triangleq
\begin{cases}
1, & \text{若 }x\in E,\\
0, & \text{否则},
\end{cases}
\end{align*}
是一个测度，称为集中于 $\{x\}$ 的 Dirac 测度。
\end{example}

% SUPPLEMENT: parent=example:ch01-计数与Dirac; kind=proof; 原书未验证 $\mu^\#$ 与 $\delta_x$ 的 $\sigma$--可加性，此处补出
\begin{proof}
(i) 由定义 $\mu^\#(\varnothing )=0$。设 $\{E_k\}\subseteq\mathcal{P}(X)$ 不交，$E\triangleq\bigcup_{k=1}^\infty E_k$。若某个 $E_k$ 为无限集，则 $E$ 为无限集，两边都等于 $\infty$。若 $E$ 为有限集，则每个 $E_k$ 有限；又因 $\{E_k\}$ 不交，$E$ 的每个元素恰属于一个 $E_k$，故至多有限多个 $E_k$ 非空，且 $\mu^\#(E)$ 等于这些非空 $E_k$ 的元素个数之和，即 $\mu^\#(E)=\sum_{k=1}^\infty\mu^\#(E_k)$。最后设 $E$ 为无限集而每个 $E_k$ 有限：此时必有无限多个 $E_k$ 非空，故右端为 $\infty=\mu^\#(E)$。总之 $\mu^\#$ 为 $\sigma$--可加的。

(ii) 由 $x\notin\varnothing$ 得 $\delta_x(\varnothing )=0$。设 $\{E_k\}\subseteq\mathcal{P}(X)$ 不交。若 $x\notin\bigcup_{k=1}^\infty E_k$，则每个 $E_k$ 都不含 $x$，两边都等于 $0$。若 $x\in\bigcup_{k=1}^\infty E_k$，则由不交性恰有一个 $k_0$ 使 $x\in E_{k_0}$，故 $\delta_x\left(\bigcup_{k=1}^\infty E_k\right)=1=\sum_{k=1}^\infty\delta_x(E_k)$。故 $\delta_x$ 为测度。
\end{proof}

% SUPPLEMENT: parent=example:ch01-计数与Dirac; kind=motivation; 说明两个测度的作用、$\sigma$--有限性的判别，以及它们与正则性、局部有限性的关系
\begin{remark}\phantomsection\label{remark:ch01-计数Dirac注}
\begin{enumerate}[(i)]
\item 引入动机。计数测度与 Dirac 测度是最基本的两个测度：前者把"元素个数"当作体积，后者把单位质量放在一点。凡涉及可数并的构造都以它们为检验对象，例如带号测度的变差（见\cref{definition:ch01-变差}）。
\item 计数测度的性质。$\mu^\#$ 有限当且仅当 $X$ 为有限集；$\mu^\#$ 为 $\sigma$--有限的充要条件是 $X$ 可数：若 $X$ 可数，取单点集 $E_k=\{x_k\}$ 即得覆盖；反之，若 $X=\bigcup_{k=1}^\infty E_k$ 且每个 $\mu^\#(E_k)<\infty$，则每个 $E_k$ 有限，故 $X$ 可数。若 $X$ 可数并枚举为 $X=\{x_1,x_2,\cdots\}$，则 $\mu^\#=\sum_{k=1}^\infty\delta_{x_k}$，即计数测度是 Dirac 测度的可数和。
\item Dirac 测度的性质。$\delta_x$ 是概率测度，且对任意 $E\subseteq X$ 有 $\delta_x(E)=\chi_E(x)$；它集中于 $\{x\}$，其支集为 $\{x\}$（见\cref{definition:ch01-Borel测度}(g)）。
\item 边界情形。在不可数的度量空间（如 $\mathbb{R}$）上，计数测度在每点的每个邻域内取值为 $\infty$，故不局部有限，从而不是\cref{definition:ch01-Borel测度}意义下的 Borel 测度；但在不可数离散空间上，它局部有限，且由\cref{remark:ch01-Borel测度性质}(ii) 内正则，因而是 Radon 测度。可见局部有限性依赖于拓扑，而不只依赖于测度本身。
\end{enumerate}
\end{remark}

\subsection{Borel 测度与 Radon 测度}\label{subsection:ch01-Borel测度}

设 $(X,\mathcal{T})$ 为一 Hausdorff 空间，$\mathcal{B}=\mathcal{B}(X)$ 为由拓扑 $\mathcal{T}$ 生成的 Borel $\sigma$--代数。以 $\mathcal{K}\subseteq\mathcal{P}(X)$ 记紧子集族。

\begin{definition}\label{definition:ch01-Borel测度}
设 $(X,\mathcal{T})$ 为一 Hausdorff 空间，$\mathcal{B}=\mathcal{B}(X)$ 为由 $\mathcal{T}$ 生成的 Borel $\sigma$--代数，$\mathcal{K}\subseteq\mathcal{P}(X)$ 为紧子集族。设 $\mathcal{A}\supseteq\mathcal{B}$ 为一 $\sigma$--代数，$\mu:\mathcal{A}\to[0,\infty]$ 为一测度。
\begin{enumerate}[(a)]
\item $\mu$ 称为局部有限的，若对任意 $x\in X$，存在 $x$ 的开邻域 $U$（从而 $U\in\mathcal{B}\subseteq\mathcal{A}$，$\mu(U)$ 有定义）使得 $\mu(U)<\infty$。局部有限的测度 $\mu:\mathcal{A}\to[0,\infty]$ 称为 $X$ 上的一个 Borel 测度。
\item 集合 $E\in\mathcal{A}$ 称为 $\mu$--外正则的（或外正则的），若
\begin{align*}
\mu(E)=\inf\{\mu(A)\mid A\supseteq E,\ A\in\mathcal{T}\}.
\end{align*}
测度 $\mu$ 称为外正则的，若每个 $E\in\mathcal{A}$ 都是 $\mu$--外正则的。
\item 集合 $E\in\mathcal{A}$ 称为 $\mu$--内正则的（或内正则的），若
\begin{align*}
\mu(E)=\sup\{\mu(K)\mid K\subseteq E,\ K\in\mathcal{K}\}.
\end{align*}
测度 $\mu$ 称为内正则的，若每个 $E\in\mathcal{A}$ 都是 $\mu$--内正则的。
\item 集合 $E\in\mathcal{A}$ 称为 $\mu$--正则的（或正则的），若它既是 $\mu$--外正则的又是 $\mu$--内正则的。测度 $\mu$ 称为正则的，若每个 $E\in\mathcal{A}$ 都是 $\mu$--正则的。
\item 内正则的 Borel 测度称为 Radon 测度。
\item Borel 测度 $\mu$ 称为适度的，若存在 $\{A_n\}\subseteq\mathcal{T}$ 使得 $X=\bigcup_{n=1}^\infty A_n$ 且对所有 $n\in\mathbb{N}$ 有 $\mu(A_n)<\infty$。
\item Borel 测度 $\mu$ 的支集（记为 $\operatorname{supp}\mu$）是满足下述条件的点 $x\in X$ 构成的集合：对 $x$ 的任意邻域 $U$ 有 $\mu(U)>0$；该集合按 $\mathcal{T}$ 是闭集。
\end{enumerate}
$X$ 上 Borel 测度与 Radon 测度的全体分别记为 $\mathcal{B}^+(X)$ 与 $\mathcal{R}^+(X)$。
\end{definition}

% SUPPLEMENT: parent=definition:ch01-Borel测度; kind=motivation; 说明 (a)--(g) 各项的动机、彼此关系（含 $\mathcal{R}^+(X)\subseteq\mathcal{B}^+(X)$）、记号约定与边界情形
\begin{remark}\phantomsection\label{remark:ch01-Borel测度动机}
\begin{enumerate}[(i)]
\item 引入动机。在拓扑空间上，测度可以用两类"简单集"来逼近：开集（从外）与紧集（从内）。正则性正是要求这种逼近对每个可测集都成立，从而使测度由其在 $\mathcal{T}$ 或 $\mathcal{K}$ 上的值完全确定。这是 Riesz 表示定理、narrow 收敛与支集理论的出发点。
\item 各项的作用。(a) 局部有限性把"测度有限"局部化到每点附近，是在局部紧空间中把 $\mu(K)<\infty$ 作为基本假设的等价形式（见\cref{remark:ch01-Borel测度性质}(iii)--(iv)）。(b) 外正则性用开集从外逼近；(c) 内正则性用紧集从内逼近；(d) 二者合起来即正则性。(e) Radon 测度 = 内正则的 Borel 测度。(f) 适度性用可数个有限测度开集覆盖 $X$，是把无穷测度分解为有限块的标准方式。(g) 支集是测度"实际占有"的闭集。
\item 由定义得 $\mathcal{R}^+(X)\subseteq\mathcal{B}^+(X)$：Radon 测度按定义为内正则的 Borel 测度。反向包含一般不成立，其成立的充分条件见\cref{proposition:ch01-可数基Borel正则}。
\item 约定与边界。(b) 中的 $\inf$ 取在非空族上（因 $X\in\mathcal{T}$ 且 $X\supseteq E$），按 $\inf\varnothing=\infty$ 的约定不影响结论；(c) 中的 $\sup$ 按 $\sup\varnothing=0$ 理解。由单调性，$\inf\{\mu(A)\mid A\supseteq E,\ A\in\mathcal{T}\}\geqslant\mu(E)$ 与 $\sup\{\mu(K)\mid K\subseteq E,\ K\in\mathcal{K}\}\leqslant\mu(E)$ 恒成立，故外（内）正则性等价于相应的反向不等式。(g) 中的集合确为闭集：其补集是那些具有 $\mu$--零测开邻域的点之集，为开集之并。(g) 也对零测度给出 $\operatorname{supp}\mu=\varnothing$，故支集可以不是 $X$；支集的进一步刻画见\cref{lemma:ch01-支集刻画}。
\end{enumerate}
\end{remark}

\begin{remark}\phantomsection\label{remark:ch01-Borel测度性质}\label{remark:ch01-局部有限}
设 $X$、$\mathcal{A}$、$\mathcal{B}$、$\mathcal{T}$、$\mathcal{K}$、$\mu$ 如\cref{definition:ch01-Borel测度}所述。
\begin{enumerate}[(i)]
\item 集合 $E\in\mathcal{A}$ 是 $\mu$--外正则的，当且仅当对任意 $\varepsilon>0$，存在 $A\in\mathcal{T}$ 使得 $A\supseteq E$ 且 $\mu(A\setminus E)<\varepsilon$。集合 $E\in\mathcal{A}$ 是 $\mu$--内正则的，当且仅当对任意 $\varepsilon>0$，存在 $K\in\mathcal{K}$ 使得 $K\subseteq E$ 且 $\mu(E\setminus K)<\varepsilon$。若 $E\in\mathcal{A}$ 是 $\mu$--外正则的，则存在 $G\in\mathcal{B}$ 使得 $G\supseteq E$ 且 $\mu(G\setminus E)=0$。
\item $\sigma$--紧集 $E\subseteq X$ 是 $\mu$--内正则的。事实上，由定义存在非减序列 $\{K_n\}\subseteq\mathcal{K}$ 使得 $E=\bigcup_{n=1}^\infty K_n$，因而 $E\in\mathcal{B}$。于是由\cref{proposition:ch01-测度性质}(iv) 有
\begin{align*}
\mu(E)=\lim_{n\to\infty}\mu(K_n)=\sup\{\mu(K)\mid K\subseteq E,\ K\in\mathcal{K}\}.
\end{align*}
\item 若 $\mu$ 局部有限（特别地，若 $\mu$ 为 Radon 测度），则对任意 $K\in\mathcal{K}$ 有 $\mu(K)<\infty$。
\item 若 Hausdorff 空间 $X$ 局部紧，则 $\mu$ 局部有限当且仅当对任意 $K\in\mathcal{K}$ 有 $\mu(K)<\infty$。特别地，若 $X$ 为局部紧 Hausdorff 空间，则 $\mu\in\mathcal{B}^+(X)$ 当且仅当对任意 $K\in\mathcal{K}$ 有 $\mu(K)<\infty$。
\item 每个适度 Borel 测度是 $\sigma$--有限的。反之，每个 $\sigma$--有限的外正则 Borel 测度是适度的。
\item 若 Hausdorff 空间 $X$ 为 $\sigma$--紧的，则 $X$ 上每个 Borel 测度 $\mu$ 都是适度的。
\end{enumerate}
\end{remark}

% SUPPLEMENT: parent=remark:ch01-Borel测度性质; kind=proof; 补出 (i) 的两种 $\varepsilon$ 判据与 $G$ 的存在性、(iii) 中开集 $A$ 的构造、(v) 的反向含义与 (vi) 的完整论证
\begin{proof}
(i) 设 $E\in\mathcal{A}$。由单调性，$\inf\{\mu(A)\mid A\supseteq E,\ A\in\mathcal{T}\}\geqslant\mu(E)$。若 $E$ 外正则且 $\mu(E)<\infty$，则对任意 $\varepsilon>0$ 存在 $A\in\mathcal{T}$，$A\supseteq E$，使 $\mu(A)<\mu(E)+\varepsilon$，于是由\cref{proposition:ch01-测度性质}(i) 与 $\mu(E)<\infty$ 得 $\mu(A\setminus E)=\mu(A)-\mu(E)<\varepsilon$。反之，若对任意 $\varepsilon>0$ 存在这样的 $A$，则 $\mu(A)<\mu(E)+\varepsilon$，故 $\inf\{\mu(A)\}=\mu(E)$，即 $E$ 外正则。当 $\mu(E)=\infty$ 时，"外正则"仅断言每个开超集的测度为 $\infty$，而条件 $\mu(A\setminus E)<\varepsilon$ 是更强的逐集要求；若 $\mu$ 为 $\sigma$--有限，取 $X=\bigcup_nX_n$、$\mu(X_n)<\infty$，对每个 $n$ 把上述论证用于 $E\cap X_n$ 得开集 $A_n\supseteq E\cap X_n$ 且 $\mu(A_n\setminus(E\cap X_n))<\varepsilon/2^n$，令 $A\triangleq\bigcup_nA_n$，则 $A\supseteq E$ 且 $\mu(A\setminus E)<\varepsilon$，故两种表述在 $\sigma$--有限情形下仍等价。

内正则性的第一种表述按同样的论证得到：由单调性 $\sup\{\mu(K)\mid K\subseteq E\}\leqslant\mu(E)$；若 $\mu(E)<\infty$ 且 $E$ 内正则，取 $K\subseteq E$ 紧使 $\mu(K)>\mu(E)-\varepsilon$，则 $\mu(E\setminus K)=\mu(E)-\mu(K)<\varepsilon$。

$G$ 的存在性：设 $E$ 外正则且 $\mu(E)<\infty$。对每个 $n\in\mathbb{N}$，由外正则性存在 $A_n\in\mathcal{T}$，$A_n\supseteq E$，使 $\mu(A_n)<\mu(E)+1/n$，于是由\cref{proposition:ch01-测度性质}(i) 与 $\mu(E)<\infty$ 得 $\mu(A_n\setminus E)<1/n$。令 $G\triangleq\bigcap_{n=1}^\infty A_n\in\mathcal{B}$。则 $G\supseteq E$，且 $G\setminus E\subseteq A_n\setminus E$ 对每个 $n$ 成立，故 $\mu(G\setminus E)\leqslant1/n$ 对每个 $n$ 成立，于是 $\mu(G\setminus E)=0$。若 $\mu(E)=\infty$ 而 $\mu$ 为 $\sigma$--有限，取 $X=\bigcup_{n=1}^\infty X_n$、$\mu(X_n)<\infty$，对每个 $n$ 把上述结论用于 $E\cap X_n$（它外正则且 $\mu$--有限），得 $G_n\in\mathcal{B}$、$G_n\supseteq E\cap X_n$、$\mu(G_n\setminus(E\cap X_n))=0$；令 $G\triangleq\bigcup_{n=1}^\infty(G_n\cap X_n)\in\mathcal{B}$，则 $G\supseteq\bigcup_{n=1}^\infty(E\cap X_n)=E$ 且 $G\setminus E\subseteq\bigcup_{n=1}^\infty(G_n\setminus(E\cap X_n))\in\mathcal{N}_\mu$，故 $\mu(G\setminus E)=0$。

(iii) 设 $K\in\mathcal{K}$。对每个 $x\in K$，由局部有限性存在开邻域 $U_x\ni x$ 使 $\mu(U_x)<\infty$。由 $K$ 的紧性，存在 $x_1,\cdots ,x_m\in K$ 使 $K\subseteq A\triangleq\bigcup_{i=1}^mU_{x_i}$。则 $A\in\mathcal{T}$、$A\supseteq K$，且由\cref{proposition:ch01-测度性质}(iii) 得 $\mu(A)\leqslant\sum_{i=1}^m\mu(U_{x_i})<\infty$；再由单调性得 $\mu(K)\leqslant\mu(A)<\infty$。

(iv) 设 $X$ 局部紧。"仅当"即 (iii)。"若"：设 $\mu(K)<\infty$ 对一切紧集 $K$ 成立。对任意 $x\in X$，由局部紧性存在 $x$ 的紧邻域 $K$，于是 $\mu(K)<\infty$，故 $\mu$ 局部有限。末句由 $\mu\in\mathcal{B}^+(X)$ 的定义即 $\mu$ 局部有限得到。

(v) 设 $\mu$ 适度：$X=\bigcup_{n=1}^\infty A_n$，$A_n\in\mathcal{T}$，$\mu(A_n)<\infty$。则 $\{A_n\}\subseteq\mathcal{A}$ 即为 $\sigma$--有限性所需的覆盖。反之，设 $\mu$ 为 $\sigma$--有限且外正则：存在 $\{E_n\}\subseteq\mathcal{A}$ 使 $X=\bigcup_{n=1}^\infty E_n$、$\mu(E_n)<\infty$。由外正则性与 (i)，对每个 $n$ 存在 $A_n\in\mathcal{T}$ 使 $A_n\supseteq E_n$ 且 $\mu(A_n\setminus E_n)<2^{-n}$，从而 $\mu(A_n)<\mu(E_n)+2^{-n}<\infty$（这里 $\mu(E_n)<\infty$，故可用\cref{proposition:ch01-测度性质}(i)）。于是 $X=\bigcup_{n=1}^\infty A_n$，$\mu$ 适度。

(vi) 设 $X=\bigcup_{n=1}^\infty K_n$，$\{K_n\}\subseteq\mathcal{K}$。由 (iii)，对每个 $n$ 存在 $A_n\in\mathcal{T}$ 使 $A_n\supseteq K_n$ 且 $\mu(A_n)<\infty$。则 $X=\bigcup_{n=1}^\infty A_n$，故 $\mu$ 适度。
\end{proof}

% SUPPLEMENT: parent=lemma:ch01-正则性; kind=motivation; 该对象的动机与条件作用见紧随其后的注记
% SUPPLEMENT: parent=lemma:ch01-紧外正则; kind=motivation; 该对象的动机与条件作用见紧随其后的注记
\begin{lemma}\label{lemma:ch01-有限Radon正则}\label{lemma:ch01-紧外正则}\label{lemma:ch01-正则性}
设 $(X,\mathcal{T})$ 为一 Hausdorff 空间，$\mu\in\mathcal{R}^+(X)$。则：
\begin{enumerate}[(i)]
\item 若 $\mu(X)<\infty$，则 $\mu$ 是正则的；
\item 每个 $K\in\mathcal{K}$ 是 $\mu$--外正则的。
\end{enumerate}
\end{lemma}

% SUPPLEMENT: parent=lemma:ch01-有限Radon正则; kind=proof; 原书证明中 $\mu(E)=\mu(X)-\mu(E^c)$、$\inf$ 与 $\sup$ 的换序、以及 (ii) 中 $\mu(M)=\mu(A)-\mu(L)$ 三处未展开，此处补全
\begin{proof}
(i) 设 $E\in\mathcal{A}$（其中 $\mathcal{A}\supseteq\mathcal{B}$），并设 $\mu(X)<\infty$。由 $\mu$ 的内正则性，$E^c\in\mathcal{A}$ 是 $\mu$--内正则的，即
\begin{align*}
\mu(E^c)=\sup\{\mu(K)\mid K\subseteq E^c,\ K\in\mathcal{K}\}.
\end{align*}
由于 $E\cup E^c=X$ 为不交并且 $\mu(X)<\infty$，由\cref{proposition:ch01-测度性质}(i) 得 $\mu(E)=\mu(X)-\mu(E^c)$。对每个紧集 $K\subseteq E^c$，其补集 $K^c\in\mathcal{T}$ 满足 $K^c\supseteq E$，且 $K\cup K^c=X$ 为不交并，故同样由\cref{proposition:ch01-测度性质}(i) 得 $\mu(K^c)=\mu(X)-\mu(K)$。于是
\begin{align*}
\mu(E)=\mu(X)-\sup_{K\subseteq E^c}\mu(K)=\inf_{K\subseteq E^c}\bigl(\mu(X)-\mu(K)\bigr)=\inf_{K\subseteq E^c}\mu(K^c),
\end{align*}
其中第二个等号用 $\mu(X)<\infty$（减去的是有限数，故上确界变为下确界）。由于
\begin{align*}
\{K^c\mid K\subseteq E^c,\ K\in\mathcal{K}\}\subseteq\{A\in\mathcal{T}\mid A\supseteq E\},
\end{align*}
由下确界的单调性得
\begin{align*}
\mu(E)\leqslant\inf\{\mu(A)\mid A\supseteq E,\ A\in\mathcal{T}\}\leqslant\inf\{\mu(K^c)\mid K\subseteq E^c,\ K\in\mathcal{K}\}=\mu(E),
\end{align*}
故 $\mu(E)=\inf\{\mu(A)\mid A\supseteq E,\ A\in\mathcal{T}\}$，即 $E$ 是 $\mu$--外正则的；结合 $\mu$ 内正则得 $\mu$ 正则。

(ii) 固定 $K\in\mathcal{K}$。由 $\mu$ 局部有限与\cref{remark:ch01-Borel测度性质}(iii)，存在 $A\in\mathcal{T}$ 使得 $A\supseteq K$ 且 $\mu(A)<\infty$。由于 $\mu$ 内正则，集合 $A\setminus K\in\mathcal{B}$ 是 $\mu$--内正则的，故对任意 $\varepsilon>0$ 存在紧集 $L\subseteq A\setminus K$ 使得
\begin{align*}
\mu(L)>\mu(A\setminus K)-\varepsilon.
\end{align*}
令 $M\triangleq A\setminus L\in\mathcal{T}$。由 $L\subseteq A\setminus K$ 得 $L\cap K=\varnothing$，故 $K\subseteq M$。又 $L\subseteq A$ 且 $\mu(A)<\infty$，由\cref{proposition:ch01-测度性质}(i) 得 $\mu(M)=\mu(A)-\mu(L)$；由 $A=K\cup(A\setminus K)$ 为不交并得 $\mu(A)=\mu(K)+\mu(A\setminus K)$。于是
\begin{align*}
\mu(M)=\mu(A)-\mu(L)<\mu(A)-\mu(A\setminus K)+\varepsilon=\mu(K)+\varepsilon,
\end{align*}
而由 $M\supseteq K$ 与\cref{proposition:ch01-测度性质}(ii) 有 $\mu(K)\leqslant\mu(M)$。故对任意 $\varepsilon>0$，
\begin{align*}
\mu(K)\leqslant\inf\{\mu(G)\mid G\supseteq K,\ G\in\mathcal{T}\}\leqslant\mu(M)<\mu(K)+\varepsilon,
\end{align*}
由 $\varepsilon$ 的任意性即得 $K$ 是 $\mu$--外正则的。
\end{proof}

% SUPPLEMENT: parent=lemma:ch01-有限Radon正则; kind=motivation; 说明两款结论的动机、有限性与局部有限性各自的作用，并指出有限性假设不可省（Dieudonné 测度）
\begin{remark}\phantomsection\label{remark:ch01-有限Radon正则注}
\begin{enumerate}[(i)]
\item 引入动机。引理把内正则性"升级"为外正则性：在有限测度情形下，内正则（紧集从内逼近）与紧集的闭集补集（开集从外逼近）互为补充，故 Radon 测度自动正则。(ii) 说明紧集本身是外正则的，这一点在构造逼近开集时反复使用。
\item 假设的作用。(i) 中 $\mu(X)<\infty$ 在两处使用：作减法 $\mu(E)=\mu(X)-\mu(E^c)$，以及把 $\sup$ 换成 $\inf$；两者都要求被减项有限。(ii) 中局部有限性只用于取得开集 $A\supseteq K$ 使 $\mu(A)<\infty$，从而使 $\mu(A)-\mu(L)$ 与 $\mu(A)-\mu(A\setminus K)$ 有意义；此处不需要 $\mu(X)<\infty$。紧集 $K$ 的闭性（Hausdorff 空间中紧集为闭集）保证 $K^c\in\mathcal{T}$，这是 (i) 中 $K^c$ 成为开超集的前提。
\item 边界情形。有限性假设不可省。在紧空间 $[0,\omega_1]$（可数序数并添最大元 $\omega_1$，赋序拓扑）上，Dieudonné 测度 $\mu$ 定义为：若 $E$ 包含 $\omega_1$ 的某个闭无界子集（相对于 $[0,\omega_1)$）则 $\mu(E)\triangleq\infty$，否则 $\mu(E)\triangleq0$。它是 Radon 测度（内正则性由"含闭无界子集"的集合包含一个紧集 $C\cup\{\omega_1\}$ 得到，局部有限性由 $[0,\alpha+1]$ 的零测性得到），且 $\mu([0,\omega_1])=\infty$。取 $E=\{\omega_1\}$，则 $\mu(E)=0$，但 $E$ 的每个开邻域都包含形如 $(\alpha,\omega_1]$ 的集合，其中含闭无界子集，故测度为 $\infty$；于是 $\inf\{\mu(A)\mid A\supseteq E,\ A\in\mathcal{T}\}=\infty\neq\mu(E)$，即 $E$ 非外正则，(i) 在 $\mu(X)=\infty$ 时失效。
\end{enumerate}
\end{remark}

\begin{proposition}\label{proposition:ch01-可数基Borel正则}
设 $(X,\mathcal{T})$ 为具有可数基的局部紧 Hausdorff 空间，$\mathcal{B}=\mathcal{B}(X)$，$\mu$ 为 $\sigma$--代数 $\mathcal{A}\supseteq\mathcal{B}$ 上的一个 Borel 测度（即 $\mu$ 局部有限）。则存在 $\{A_n\}\subseteq\mathcal{T}$ 使得 $X=\bigcup_{n=1}^\infty A_n$ 且 $\mu(A_n)<\infty$（$n\in\mathbb{N}$），且每个 $E\in\mathcal{A}$ 既是 $\mu$--外正则的又是 $\mu$--内正则的。
\end{proposition}

% SUPPLEMENT: parent=proposition:ch01-可数基Borel正则; kind=proof; 记录证明中两次使用的附录 A 事实（可数基局部紧空间中开子集 $\sigma$--紧）
\begin{remark}\phantomsection\label{remark:ch01-附录A1}
由\cref{remark:appA-可数基局部紧}，具有可数基的局部紧 Hausdorff 空间中每个开子集都是 $\sigma$--紧的；取该开子集为 $X$ 本身即知 $X$ 是 $\sigma$--紧的。该事实在下述证明中使用两次：第一次用于由 $\sigma$--紧性得到适度性，第二次用于断言每个 $A_n$ 的开子集 $\sigma$--紧，从而 $\mu_n$ 内正则。因此可数基假设不是技术性的：它保证拓扑由可数个集合生成，从而把 $\sigma$--紧性传递给一切开子集；缺少它时（例如在\cref{remark:ch01-有限Radon正则注}所述的 $[0,\omega_1)$ 上）论证中断，正则性确实可能失效。
\end{remark}

% SUPPLEMENT: parent=proposition:ch01-可数基Borel正则; kind=proof; 原书证明中"如\cref{lemma:ch01-有限Radon正则}(i) 那样论证"与"$A_n'\in\mathcal{T}$"两处为跳跃，此处补出子空间上的有限 Radon 测度论证与 $A_n'$ 的开性
\begin{proof}
由\cref{remark:appA-可数基局部紧}，$X$ 是 $\sigma$--紧的，故由\cref{remark:ch01-Borel测度性质}(vi)，$\mu$ 是适度的。

为证 $\mu$ 外正则，取 $\{A_n\}\subseteq\mathcal{T}$，使 $X=\bigcup_{n=1}^\infty A_n$ 且 $\mu(A_n)<\infty$（不妨设 $\{A_n\}$ 非减，否则用 $A_1\cup\cdots\cup A_n$ 代替 $A_n$）。设 $E\in\mathcal{A}\supseteq\mathcal{B}$，令 $E_n\triangleq E\cap A_n$。则 $E=\bigcup_{n=1}^\infty E_n$ 且 $\mu(E_n)\leqslant\mu(A_n)<\infty$。

固定 $n$。由\cref{remark:ch01-测度子空间}，$\mu_n\triangleq\mu|_{\mathcal{B}(A_n)}$ 是迹 $\sigma$--代数 $\mathcal{B}(A_n)=\mathcal{B}\cap A_n$ 上的测度，且 $\mu_n(A_n)=\mu(A_n)<\infty$。在 Hausdorff 空间 $A_n$（赋子空间拓扑）上，$\mu_n$ 是有限 Borel 测度；由\cref{remark:appA-可数基局部紧}，$A_n$ 的每个开子集是 $\sigma$--紧的，故由\cref{remark:ch01-Borel测度性质}(ii) 是 $\mu_n$--内正则的，再由 $\mu_n$ 的 $\sigma$--可加性知 $\mu_n$ 内正则。对空间 $A_n$、测度 $\mu_n$ 重复\cref{lemma:ch01-有限Radon正则}(i) 的论证（该论证只用到有限性与内正则性），可知 $\mu_n$ 正则，特别地外正则：对任意 $\varepsilon>0$ 存在 $A_n'\in\mathcal{T}$ 使得
\begin{align*}
E_n\subseteq A_n'\subseteq A_n,\qquad\mu_n(A_n'\setminus E_n)=\mu(A_n'\setminus E_n)<\frac{\varepsilon}{2^n}.
\end{align*}
此处 $A_n'$ 是 $A_n$ 中的开集即可，而 $A_n\in\mathcal{T}$，故 $A_n'$ 是 $A_n$ 与某个开集之交，从而 $A_n'\in\mathcal{T}$。

令 $A'\triangleq\bigcup_{n=1}^\infty A_n'\in\mathcal{T}$。则 $A'\supseteq\bigcup_{n=1}^\infty E_n=E$，且由\cref{proposition:ch01-测度性质}(iii)
\begin{align*}
\mu(A'\setminus E)\leqslant\sum_{n=1}^\infty\mu(A_n'\setminus E_n)<\varepsilon.
\end{align*}
故 $\mu$ 是外正则的（见\cref{remark:ch01-Borel测度性质}(i)）。

现证 $\mu$ 内正则。设 $E\in\mathcal{A}$，$E_n\triangleq E\cap A_n$，$\mu_n$ 同上。对任意 $\alpha<\mu(E)$，令 $B_p\triangleq\bigcup_{n=1}^pE_n$（$p\in\mathbb{N}$）。序列 $\{B_p\}$ 非减且 $\bigcup_{p=1}^\infty B_p=E$，故由\cref{proposition:ch01-测度性质}(iv) 得 $\mu(B_p)\to\mu(E)$；于是可取 $p$ 充分大使得 $\mu(B_p)>\alpha$，并令 $\varepsilon\triangleq\mu(B_p)-\alpha>0$（注意 $\mu(B_p)<\infty$，因为 $B_p\subseteq\bigcup_{n=1}^pA_n$）。对每个 $n=1,\cdots ,p$，测度 $\mu_n$ 正则，故内正则，于是存在紧集 $K_n\subseteq E_n$ 使得
\begin{align*}
\mu_n(E_n\setminus K_n)=\mu(E_n\setminus K_n)<\frac{\varepsilon}{p}.
\end{align*}
令 $K\triangleq\bigcup_{n=1}^pK_n$。则 $K\in\mathcal{K}$（有限个紧集之并）、$K\subseteq B_p$，且由\cref{proposition:ch01-测度性质}(iii)
\begin{align*}
\mu\left(B_p\setminus K\right)\leqslant\sum_{n=1}^p\mu(E_n\setminus K_n)<\varepsilon=\mu(B_p)-\alpha.
\end{align*}
由 $B_p=K\cup(B_p\setminus K)$ 为不交并与\cref{proposition:ch01-测度性质}(i) 得 $\mu(K)=\mu(B_p)-\mu(B_p\setminus K)>\alpha$。于是对任意 $\alpha<\mu(E)$ 存在紧集 $K\subseteq E$ 使 $\mu(K)>\alpha$，故 $\sup\{\mu(K)\mid K\subseteq E,\ K\in\mathcal{K}\}\geqslant\mu(E)$；反向不等式由单调性成立，故 $E$ 内正则。
\end{proof}

% SUPPLEMENT: parent=proposition:ch01-可数基Borel正则; kind=motivation; 说明可数基假设与局部紧性的作用、结论中"正则"与"适度"的关系，以及缺少可数基时的反例
\begin{remark}\phantomsection\label{remark:ch01-可数基Borel正则注}
\begin{enumerate}[(i)]
\item 引入动机。命题说明在"足够好"的空间（$\mathbb{R}^N$、可数基局部紧 Hausdorff 空间）中，任一 Borel 测度自动正则且适度，因此不必逐例验证正则性。它的结论在\cref{remark:ch01-BorelRadon一致}中化为"Borel 测度与 Radon 测度重合"。
\item 假设的作用。局部紧 Hausdorff 保证紧集足够多且每点的紧邻域存在；可数基经\cref{remark:appA-可数基局部紧}提供 $\sigma$--紧性与"开子集仍 $\sigma$--紧"这一继承性质，后者使\cref{remark:ch01-Borel测度性质}(ii) 可应用于每个 $A_n$ 的开子集。适度性则把 $X$ 分解为可数个有限测度块，使\cref{lemma:ch01-有限Radon正则}可逐块使用；两处使用缺一不可。
\item 边界情形。可数基假设不可省。在 $[0,\omega_1)$（第一不可数序数空间，局部紧 Hausdorff 但无可数基）上，Dieudonné 测度（含闭无界子集者取 $\infty$、其余取 $0$）是局部有限的 Borel 测度，但 $\mu([0,\omega_1))=\infty$，而每个紧集有界故测度为 $0$，于是内正则性在 $E=[0,\omega_1)$ 处失效，$\mu$ 不是 Radon 测度，更不正则。
\end{enumerate}
\end{remark}

\begin{remark}\phantomsection\label{remark:ch01-BorelRadon一致}
由\cref{proposition:ch01-可数基Borel正则}，在具有可数基的局部紧 Hausdorff 空间（例如 $\mathbb{R}^N$）中，Borel 测度与 Radon 测度重合，且都是正则的。
\end{remark}

% SUPPLEMENT: parent=remark:ch01-BorelRadon一致; kind=proof; 原书只陈述重合性，此处补出 $\mathcal{B}^+(X)=\mathcal{R}^+(X)$ 的两个包含
\begin{proof}
设 $X$ 为具有可数基的局部紧 Hausdorff 空间。由\cref{definition:ch01-Borel测度}(e)，每个 Radon 测度按定义是内正则的 Borel 测度，故 $\mathcal{R}^+(X)\subseteq\mathcal{B}^+(X)$。反之，设 $\mu\in\mathcal{B}^+(X)$。由\cref{proposition:ch01-可数基Borel正则}，$\mu$ 正则，特别地内正则；又 $\mu$ 按定义为 Borel 测度，故 $\mu$ 是内正则的 Borel 测度，即 $\mu\in\mathcal{R}^+(X)$。于是 $\mathcal{B}^+(X)=\mathcal{R}^+(X)$，且由同一命题两者都正则。
\end{proof}

\subsection{零测集}\label{subsection:ch01-零测集}

\begin{definition}\label{definition:ch01-零测集}
设 $(X,\mathcal{A},\mu)$ 为一测度空间。
\begin{enumerate}[(i)]
\item 集合 $N\subseteq X$ 称为 $\mu$--零测的（或零测的），若 $N\in\mathcal{A}$ 且 $\mu(N)=0$。集合 $E\subseteq X$ 称为 $\mu$--可忽略的（或可忽略的），若存在 $\mu$--零测集 $N$ 使得 $E\subseteq N$。
\item 空间 $(X,\mathcal{A},\mu)$ 称为完备的，若所有可忽略集都是可测的，因而是零测的。此时称 $\mu$ 为完备测度，$\sigma$--代数 $\mathcal{A}$ 关于 $\mu$ 是完备的。
\end{enumerate}
\end{definition}

记零测集族为 $\mathcal{N}_\mu$，可忽略集族为 $\mathcal{U}_\mu$。两族 $\mathcal{N}_\mu$ 与 $\mathcal{U}_\mu$ 关于可数并都是稳定的。此外，$\mathcal{U}_\mu$ 是遗传的，即 $F\subseteq E$ 且 $E\in\mathcal{U}_\mu$ 蕴含 $F\in\mathcal{U}_\mu$。又 $\mathcal{N}_\mu=\mathcal{U}_\mu$ 当且仅当 $(X,\mathcal{A},\mu)$ 是完备的。

% SUPPLEMENT: parent=definition:ch01-零测集; kind=proof; 补出上段四条断言：$\mathcal{N}_\mu$ 与 $\mathcal{U}_\mu$ 对可数并稳定、$\mathcal{U}_\mu$ 的遗传性、以及 $\mathcal{N}_\mu=\mathcal{U}_\mu$ 与完备性的等价
\begin{proof}
$\mathcal{N}_\mu$ 对可数并稳定：设 $\{N_k\}\subseteq\mathcal{N}_\mu$。则 $\bigcup_{k=1}^\infty N_k\in\mathcal{A}$ 且由\cref{proposition:ch01-测度性质}(iii) 有 $\mu\left(\bigcup_{k=1}^\infty N_k\right)\leqslant\sum_{k=1}^\infty\mu(N_k)=0$。

$\mathcal{U}_\mu$ 对可数并稳定：设 $E_k\subseteq N_k$，$N_k\in\mathcal{N}_\mu$。则 $N\triangleq\bigcup_{k=1}^\infty N_k\in\mathcal{N}_\mu$ 且 $\bigcup_{k=1}^\infty E_k\subseteq N$，故 $\bigcup_{k=1}^\infty E_k\in\mathcal{U}_\mu$。

遗传性：若 $F\subseteq E\subseteq N$ 且 $N\in\mathcal{N}_\mu$，则 $F\subseteq N$，故 $F\in\mathcal{U}_\mu$。

等价性：若 $\mathcal{N}_\mu=\mathcal{U}_\mu$，则每个可忽略集属于 $\mathcal{N}_\mu$，即可测且零测，故 $(X,\mathcal{A},\mu)$ 完备。反之，若空间完备，则 $\mathcal{U}_\mu\subseteq\mathcal{N}_\mu$；而 $\mathcal{N}_\mu\subseteq\mathcal{U}_\mu$ 由定义成立（取 $N=E$）。故 $\mathcal{N}_\mu=\mathcal{U}_\mu$。
\end{proof}

% SUPPLEMENT: parent=definition:ch01-零测集; kind=motivation; 说明零测与可忽略之别、完备化的动机，并给出可忽略集不可测的例子
\begin{remark}\phantomsection\label{remark:ch01-零测集注}
\begin{enumerate}[(i)]
\item 引入动机。测度为零的集合不影响积分值，因此"例外集零测"是分析中讨论"几乎处处"性质的基础。定义把两类集合分开：零测集必须可测，可忽略集只是零测集的子集，未必可测。这一区分在非完备空间中是真区别（见 (iii)），也正是引入完备化的动机。
\item 假设的作用。"$N\in\mathcal{A}$"是零测性的组成部分：若只要求 $\mu(N)=0$ 而不要求可测，则 $\mu(N)$ 本身没有定义。完备性要求"可忽略 $\Rightarrow$ 可测"，等价于 $\mathcal{N}_\mu=\mathcal{U}_\mu$（见上方证明），从而在完备空间中两种说法可以互换使用。
\item 边界情形。取 $X=\{0,1\}$、$\mathcal{A}=\{\varnothing ,X\}$、$\mu(\varnothing )=0$、$\mu(X)=1$。则 $N=\{0\}$ 是可忽略集（包含于零测集 $\varnothing$），但 $N\notin\mathcal{A}$，故空间不完备，且 $\mathcal{U}_\mu\neq\mathcal{N}_\mu$。若改取 $\mu\equiv0$，则 $\mathcal{U}_\mu=\mathcal{P}(X)$ 而 $\mathcal{N}_\mu=\mathcal{A}$，两者相等当且仅当 $\mathcal{A}=\mathcal{P}(X)$。
\end{enumerate}
\end{remark}

\begin{theorem}\label{theorem:ch01-Lebesgue完备化}
设 $(X,\mathcal{A},\mu)$ 为一测度空间。则：
\begin{enumerate}[(i)]
\item 族
\begin{align}
\overline{\mathcal{A}}\triangleq \{E\subseteq X\mid\exists F,G\in\mathcal{A}\text{ 使得 }F\subseteq E\subseteq G,\ \mu(G\setminus F)=0\} \label{align:ch01-完备sigma代数}
\end{align}
是包含 $\mathcal{A}$ 的 $\sigma$--代数；
\item 存在完备测度 $\overline{\mu}:\overline{\mathcal{A}}\to[0,\infty]$ 使得 $\overline{\mu}|_{\mathcal{A}}=\mu$。
\end{enumerate}
此外，若 $(X,\mathcal{A}_1,\mu_1)$ 为完备测度空间且 $\mathcal{A}_1\supseteq\mathcal{A}$、$\mu_1|_{\mathcal{A}}=\mu$，则 $\mathcal{A}_1\supseteq\overline{\mathcal{A}}$ 且 $\mu_1|_{\overline{\mathcal{A}}}=\overline{\mu}$；即 $(X,\overline{\mathcal{A}},\overline{\mu})$ 是包含 $(X,\mathcal{A},\mu)$ 的最小完备测度空间。
\end{theorem}

% SUPPLEMENT: parent=definition:ch01-Lebesgue完备化; kind=motivation; 该对象的动机与条件作用见紧随其后的注记
% SUPPLEMENT: parent=definition:ch01-完备化; kind=motivation; 该对象的动机与条件作用见紧随其后的注记
\begin{definition}\label{definition:ch01-完备化}\label{definition:ch01-Lebesgue完备化}
设 $(X,\mathcal{A},\mu)$ 为一测度空间，$\overline{\mathcal{A}}$ 与 $\overline{\mu}$ 由\cref{theorem:ch01-Lebesgue完备化} 给出。测度空间 $(X,\overline{\mathcal{A}},\overline{\mu})$ 称为 $(X,\mathcal{A},\mu)$ 的 Lebesgue 完备化。由该定理的极小性，$\overline{\mu}$ 是 $\mu$ 在 $\mathcal{A}$ 之外唯一的完备延拓，故 $\overline{\mu}$ 也简称为 $\mu$ 的 Lebesgue 完备化。
\end{definition}

% SUPPLEMENT: parent=theorem:ch01-Lebesgue完备化; kind=proof; 原证明中"由 (1.2) 立即得 $E^c\in\overline{\mathcal{A}}$"、$\overline{\mu}$ 的 $\sigma$--可加性所依赖的 $\bigcup_nG_n\setminus\bigcup_nF_n\subseteq\bigcup_n(G_n\setminus F_n)$、以及末尾"其余结论显然"处 $\mu_1|_{\overline{\mathcal{A}}}=\overline{\mu}$ 的论证均未展开，此处补全
\begin{proof}
(i) 由 $\mathcal{A}\subseteq\overline{\mathcal{A}}$（取 $F=G=E$）知 $\mathcal{A}$ 含于 $\overline{\mathcal{A}}$，特别地 $\varnothing \in\overline{\mathcal{A}}$。

对补集封闭：设 $E\in\overline{\mathcal{A}}$，取 $F,G\in\mathcal{A}$ 使 $F\subseteq E\subseteq G$、$\mu(G\setminus F)=0$。则 $G^c\subseteq E^c\subseteq F^c$，$G^c,F^c\in\mathcal{A}$，且 $F^c\setminus G^c=G\setminus F$，故 $\mu(F^c\setminus G^c)=0$，即 $E^c\in\overline{\mathcal{A}}$。

对可数并封闭：设 $\{E_n\}\subseteq\overline{\mathcal{A}}$，取 $\{F_n\},\{G_n\}\subseteq\mathcal{A}$ 使 $F_n\subseteq E_n\subseteq G_n$ 且 $\mu(G_n\setminus F_n)=0$。令 $F\triangleq\bigcup_{n=1}^\infty F_n$、$G\triangleq\bigcup_{n=1}^\infty G_n$。则 $F,G\in\mathcal{A}$ 且 $F\subseteq\bigcup_{n=1}^\infty E_n\subseteq G$。又若 $x\in G\setminus F$，则 $x\in G_n$ 对某个 $n$ 成立，而 $x\notin F$ 给出 $x\notin F_n$，故 $x\in G_n\setminus F_n$；于是
\begin{align*}
G\setminus F\subseteq\bigcup_{n=1}^\infty(G_n\setminus F_n),
\end{align*}
由 $\mu$ 的 $\sigma$--次可加性得 $\mu(G\setminus F)\leqslant\sum_{n=1}^\infty\mu(G_n\setminus F_n)=0$。故 $\bigcup_{n=1}^\infty E_n\in\overline{\mathcal{A}}$。于是 $\overline{\mathcal{A}}$ 是包含 $\mathcal{A}$ 的 $\sigma$--代数。

(ii) 对 $E\in\overline{\mathcal{A}}$，取 $F,G\in\mathcal{A}$ 使 $F\subseteq E\subseteq G$、$\mu(G\setminus F)=0$，并令 $\overline{\mu}(E)\triangleq\mu(F)$。先注意 $\mu(F)=\mu(G)$：由 $G=F\cup(G\setminus F)$ 为不交并及\cref{proposition:ch01-测度性质}(i)，有 $\mu(G)=\mu(F)+\mu(G\setminus F)=\mu(F)$。故 $\overline{\mu}(E)$ 也可写成 $\mu(G)$。

良定义性：设 $F_1,G_1\in\mathcal{A}$ 亦满足 $F_1\subseteq E\subseteq G_1$、$\mu(G_1\setminus F_1)=0$。则 $F_1\subseteq E\subseteq G$ 且 $F\subseteq E\subseteq G_1$；若 $x\in G\setminus F_1$，则或者 $x\notin F$（故 $x\in G\setminus F$），或者 $x\in F\subseteq G_1$（故 $x\in G_1\setminus F_1$），于是
\begin{align*}
G\setminus F_1\subseteq(G\setminus F)\cup(G_1\setminus F_1).
\end{align*}
由此得 $\mu(G\setminus F_1)=0$，同法得 $\mu(G_1\setminus F)=0$。于是由 $\mu$ 的有限可加性
\begin{align*}
\mu(G)=\mu(F_1)+\mu(G\setminus F_1)=\mu(F_1),\qquad\mu(G_1)=\mu(F)+\mu(G_1\setminus F)=\mu(F),
\end{align*}
故 $\mu(F)=\mu(F_1)$，即 $\overline{\mu}(E)$ 不依赖于 $F,G$ 的选取。

$\overline{\mu}$ 是测度：由 $\varnothing \subseteq\varnothing \subseteq\varnothing$ 得 $\overline{\mu}(\varnothing )=\mu(\varnothing )=0$。设 $\{E_n\}\subseteq\overline{\mathcal{A}}$ 不交，取 $F_n,G_n$ 如上。则 $\{F_n\}\subseteq\mathcal{A}$ 不交（因 $F_n\subseteq E_n$），且由 (i) 中已证的事实，$F\triangleq\bigcup_nF_n\subseteq\bigcup_nE_n\subseteq\bigcup_nG_n\triangleq G$ 满足 $\mu(G\setminus F)=0$，故
\begin{align*}
\overline{\mu}\left(\bigcup_{n=1}^\infty E_n\right)=\mu(F)=\mu\left(\bigcup_{n=1}^\infty F_n\right)=\sum_{n=1}^\infty\mu(F_n)=\sum_{n=1}^\infty\overline{\mu}(E_n),
\end{align*}
其中第三个等号用 $\mu$ 的 $\sigma$--可加性。

完备性：设 $N\in\overline{\mathcal{A}}$、$\overline{\mu}(N)=0$，并设 $E\subseteq N$。取 $F,G\in\mathcal{A}$ 使 $F\subseteq N\subseteq G$、$\mu(G\setminus F)=0$。由 $\overline{\mu}(N)=\mu(G)$ 得 $\mu(G)=0$。于是 $\varnothing \subseteq E\subseteq G$ 且 $\mu(G\setminus\varnothing )=\mu(G)=0$，故 $E\in\overline{\mathcal{A}}$，即 $\overline{\mu}$ 完备。

极小性：设 $(X,\mathcal{A}_1,\mu_1)$ 为完备测度空间，$\mathcal{A}_1\supseteq\mathcal{A}$，$\mu_1|_{\mathcal{A}}=\mu$。设 $E\in\overline{\mathcal{A}}$，取 $F,G\in\mathcal{A}$ 使 $F\subseteq E\subseteq G$、$\mu(G\setminus F)=0$。由 $\mu_1$ 完备、$\mu_1(G\setminus F)=\mu(G\setminus F)=0$ 与 $G\setminus E\subseteq G\setminus F$ 得 $G\setminus E\in\mathcal{A}_1$；于是 $E=G\setminus(G\setminus E)\in\mathcal{A}_1$，故 $\overline{\mathcal{A}}\subseteq\mathcal{A}_1$。

最后证 $\mu_1|_{\overline{\mathcal{A}}}=\overline{\mu}$。设 $E\in\overline{\mathcal{A}}$，$F,G\in\mathcal{A}$ 如\eqref{align:ch01-完备sigma代数}。由 $G\setminus E\subseteq G\setminus F$ 与 $\mu_1(G\setminus F)=0$ 得 $\mu_1(G\setminus E)=0$；又 $G=E\cup(G\setminus E)$ 为不交并且 $E,G\setminus E\in\mathcal{A}_1$，故由 $\mu_1$ 的有限可加性 $\mu_1(G)=\mu_1(E)+\mu_1(G\setminus E)=\mu_1(E)$。于是
\begin{align*}
\mu_1(E)=\mu_1(G)=\mu(G)=\overline{\mu}(E),
\end{align*}
其中第二个等号用 $\mu_1|_{\mathcal{A}}=\mu$。故 $\mu_1|_{\overline{\mathcal{A}}}=\overline{\mu}$。
\end{proof}

% SUPPLEMENT: parent=theorem:ch01-Lebesgue完备化; kind=motivation; 说明完备化的动机（零测集的子集应可测）、两处关键构造的作用、极小性的含义与边界
\begin{remark}\phantomsection\label{remark:ch01-完备化动机}
\begin{enumerate}[(i)]
\item 引入动机。非完备空间中，零测集的子集可能不可测，而分析中经常需要"忽略零测集"：例如两个 $\mu$--a.e. 相等的函数之差只在可忽略集上非零，只有在该集可测时才能断言其积分为零。完备化以最小的代价把所有此类子集纳入 $\sigma$--代数，而不改变原有集合上的测度值。
\item 构造的作用。条件"$F\subseteq E\subseteq G$ 且 $\mu(G\setminus F)=0$"同时表达了两件事：$E$ 与可测集 $F,G$ 只差一个零测集，因而 $E$ 应当可测；$\overline{\mu}(E)\triangleq\mu(F)$ 则说明 $E$ 的测度由 $F$ 与 $G$ 唯一确定（这两个值相等，原因见上方证明）。极小性表明：(i) 中的 $\sigma$--代数已是容纳一切可忽略集的最小可能选择，而 $\overline{\mu}$ 由 $\mu$ 唯一确定，故完备化是典范的。
\item 边界情形。完备化可能严格扩大 $\sigma$--代数：取 $X=\{0,1\}$、$\mathcal{A}=\{\varnothing ,X\}$、$\mu\equiv0$，则对 $E=\{0\}$ 有 $\varnothing \subseteq E\subseteq X$ 且 $\mu(X\setminus\varnothing )=0$，故 $E\in\overline{\mathcal{A}}$，于是 $\overline{\mathcal{A}}=\mathcal{P}(X)\supsetneq\mathcal{A}$。若把上例改为 $\mu(X)=1$（其余不变），则 $\{0\}\notin\overline{\mathcal{A}}$：满足 $F\subseteq\{0\}\subseteq G$ 的 $F,G\in\mathcal{A}$ 只能是 $F=\varnothing$、$G=X$，而 $\mu(X\setminus\varnothing )=1\neq0$，故此时 $\overline{\mathcal{A}}=\mathcal{A}$。可见完备化是否扩大 $\sigma$--代数取决于 $\mu$ 的零测集。最后，若 $\mu$ 本身完备，则对 $E\in\overline{\mathcal{A}}$ 与相应的 $F,G$ 有 $E\setminus F\subseteq G\setminus F\in\mathcal{N}_\mu$，由完备性得 $E\setminus F\in\mathcal{A}$，故 $E=F\cup(E\setminus F)\in\mathcal{A}$，即 $\overline{\mathcal{A}}=\mathcal{A}$ 且 $\overline{\mu}=\mu$。
\end{enumerate}
\end{remark}

\begin{remark}\phantomsection\label{remark:ch01-完备化刻画}
\eqref{align:ch01-完备sigma代数} 中定义的 $\sigma$--代数 $\overline{\mathcal{A}}$ 有如下刻画：
\begin{align*}
\overline{\mathcal{A}}=\{F\cup E_0\mid F\in\mathcal{A},\ E_0\in\mathcal{U}_\mu\}.
\end{align*}
事实上，设 $E\in\overline{\mathcal{A}}$，并取 $F,G\in\mathcal{A}$ 使 $F\subseteq E\subseteq G$、$\mu(G\setminus F)=0$。则 $E_0\triangleq E\setminus F\subseteq G\setminus F$，而 $G\setminus F\in\mathcal{N}_\mu$，故 $E_0\in\mathcal{U}_\mu$ 且 $E=F\cup E_0$。反之，设 $E=F\cup E_0$，其中 $F\in\mathcal{A}$、$E_0\in\mathcal{U}_\mu$。由 $\mathcal{U}_\mu$ 的定义存在 $N_0\in\mathcal{N}_\mu$ 使 $E_0\subseteq N_0$；令 $G\triangleq F\cup(N_0\setminus F)\in\mathcal{A}$。则 $F\subseteq E\subseteq G$ 且 $\mu(G\setminus F)=\mu(N_0\setminus F)=0$，故 $E\in\overline{\mathcal{A}}$。
\end{remark}

\begin{definition}\label{definition:ch01-a.e.}
设 $(X,\mathcal{A},\mu)$ 为一测度空间。性质 $P$（对 $x\in X$ 陈述）称为 $\mu$--几乎处处成立（记为 $\mu$--a.e. 或 a.e.），若集合 $\{x\in X\mid P(x)\text{ 不成立}\}$ 属于 $\mathcal{N}_\mu$；即例外集既 $\mu$--可测又 $\mu$--零测。
\end{definition}

作为上述定义的特殊情形：(a) 两个函数 $f,g:X\to\mathbb{R}$ 在 $X$ 上 $\mu$--a.e. 相等，若 $\{x\in X\mid f(x)\neq g(x)\}\in\mathcal{N}_\mu$；(b) 序列 $\{f_n\}$（$f_n:X\to\mathbb{R}$）在 $X$ 中 $\mu$--a.e. 收敛，若 $\{x\in X\mid\{f_n(x)\}\text{ 不收敛}\}\in\mathcal{N}_\mu$。

$\mu$--a.e. 相等是一个等价关系：
\begin{align}
g\sim f\overset{\mathrm{def}}{\Longleftrightarrow} g=f\ \mu\text{--a.e. 在 }X\text{ 中}. \label{align:ch01-a.e.等价}
\end{align}
关于\eqref{align:ch01-a.e.等价} 的每个等价类由其任意一个元素唯一确定，该元素称为该类的一个代表元。因此自然地把整个类视为一个在 $X$ 中 $\mu$--a.e. 定义的唯一映射 $f$。

% SUPPLEMENT: parent=definition:ch01-a.e.; kind=motivation; 说明"例外集属于 $\mathcal{N}_\mu$"这一定法的用意（与可忽略集版本的差别）、与积分理论的联系，以及边界情形
\begin{remark}\phantomsection\label{remark:ch01-a.e.动机}
\begin{enumerate}[(i)]
\item 引入动机。把"几乎处处"形式化，才能对 a.e. 成立的命题作逻辑运算，并讨论 a.e. 相等的等价类（$L^p$ 空间的元素正是这样的类）。定义要求例外集属于 $\mathcal{N}_\mu$，即例外集既可测又零测。
\item 假设的作用。要求例外集属于 $\mathcal{N}_\mu$ 而非仅仅"包含于某个零测集"，使定义在非完备空间中仍然是可检验的可测性陈述：若只要求例外集可忽略，则在空间不完备时无法由 $\mu$ 直接验证。完备化（\cref{theorem:ch01-Lebesgue完备化}）之后两种表述一致，因为此时 $\mathcal{N}_\mu=\mathcal{U}_\mu$。
\item 边界情形。若 $\mu\equiv0$，则 $\mathcal{N}_\mu$ 由 $\mathcal{A}$ 中一切集合组成，故任何性质都 $\mu$--a.e. 成立；反之，若 $P$ 在某个非空开集上处处不成立而 $\mu$ 在该开集上为正，则 $P$ 不 a.e. 成立。$\mu$--a.e. 相等的两支函数之差不必处处为零，故"$f=g$ a.e."不是逐点等式；这正是在 $L^p$ 理论中把 a.e. 相等的函数等同起来的原因。
\end{enumerate}
\end{remark}

% SUPPLEMENT: parent=definition:ch01-a.e.; kind=proof; 原书称 $\mu$--a.e. 相等为等价关系而未验证，此处补出自反、对称、传递三条
\begin{proof}
对函数 $f,g,h:X\to\mathbb{R}$，记 $N(f,g)\triangleq\{x\in X\mid f(x)\neq g(x)\}$。

自反性：$N(f,f)=\varnothing\in\mathcal{N}_\mu$，故 $f\sim f$。对称性：$N(g,f)=N(f,g)\in\mathcal{N}_\mu$，故 $g\sim f$ 当且仅当 $f\sim g$。

传递性：设 $f\sim g$、$g\sim h$。若 $x\notin N(f,g)\cup N(g,h)$，则 $f(x)=g(x)=h(x)$，故 $N(f,h)\subseteq N(f,g)\cup N(g,h)$。由 $\mathcal{N}_\mu$ 对可数并稳定（见\cref{definition:ch01-零测集}后的证明）得 $N(f,g)\cup N(g,h)\in\mathcal{N}_\mu$，故 $N(f,h)\in\mathcal{N}_\mu$，即 $f\sim h$。
\end{proof}

\begin{definition}\label{definition:ch01-a.e.定义}
设 $(X,\mathcal{A},\mu)$ 为一测度空间，$Y$ 为一集合。函数 $f:E\to Y$（$E\in\mathcal{A}$）称为在 $X$ 中 $\mu$--a.e. 定义的，若 $E^c\in\mathcal{N}_\mu$；即 $f$ 的定义域可测，且其补集为 $\mu$--零测集。
\end{definition}

% SUPPLEMENT: parent=definition:ch01-a.e.定义; kind=motivation; 说明"定义域补集零测"的含义、与 a.e. 相等的关系、古典代表的用法与边界
\begin{remark}\phantomsection\label{remark:ch01-a.e.定义动机}
\begin{enumerate}[(i)]
\item 引入动机。把 $\mu$--a.e. 相等的函数等同为一个对象之后，这个对象自然只在 $X$ 的一个"满测度"子集上有定义，定义\cref{definition:ch01-a.e.定义} 正是为此类对象提供载体：$f$ 的定义域 $E$ 可测，且其补集零测。
\item 假设的作用。要求 $E\in\mathcal{A}$ 与 $E^c\in\mathcal{N}_\mu$ 使 $f$ 与 $\mu$ 相容：只在可测集上取值，且例外集可测零测。这保证了"$f$ 在 $X$ 中 a.e. 有定义"这一断言属于 $\mathcal{A}$ 可检验的范围。$Y$ 为任意集合，故该定义也适用于取值于度量空间或 Banach 空间的函数。
\item 边界情形。$E=X$ 时 $E^c=\varnothing\in\mathcal{N}_\mu$，故每个经典函数都是 a.e. 定义的；但 a.e. 定义的函数未必有经典代表（当 $\mu$ 不完备且 $E$ 之外无法延拓时）。若 $\mu(X)=0$，则任何定义在可测集上的函数都 a.e. 有定义。由于 $E\in\mathcal{A}$ 蕴含 $E^c\in\mathcal{A}$，这里的条件 $E^c\in\mathcal{N}_\mu$ 与 $\mu(E^c)=0$ 等价。
\end{enumerate}
\end{remark}

在适当场合，为强调\cref{definition:ch01-a.e.定义} 与经典逐点定义之间的区别，我们说：
\begin{enumerate}[(i)]
\item 函数 $f:X\to\mathbb{R}$ 是经典的，若它在每个点 $x\in X$ 处都有定义；
\item 经典函数 $g$ 等于一个 $\mu$--a.e. 定义的函数 $f$，若 $g$ 是 $f$ 的一个经典代表元。
\end{enumerate}

在此联系下，序列的 $\mu$--a.e. 极限在\cref{definition:ch01-a.e.定义} 的意义下是唯一的（但在逐点意义下不唯一）。

% SUPPLEMENT: parent=definition:ch01-a.e.定义; kind=proof; 补出"a.e. 极限唯一"：$f_n\to f$ a.e. 与 $f_n\to g$ a.e. 蕴含 $f=g$ a.e.
\begin{proof}
设 $N_1\triangleq\{x\in X\mid\{f_n(x)\}\text{ 不收敛}\}\in\mathcal{N}_\mu$，并设 $f_n\to f$ $\mu$--a.e.、$f_n\to g$ $\mu$--a.e.。又设 $f,g$ 分别在 $E_f,E_g\in\mathcal{A}$ 上有定义且 $E_f^c,E_g^c\in\mathcal{N}_\mu$。

令 $N\triangleq N_1\cup E_f^c\cup E_g^c$。由 $\mathcal{N}_\mu$ 对可数并稳定得 $N\in\mathcal{N}_\mu$。若 $x\notin N$，则 $x\in E_f\cap E_g$、$\{f_n(x)\}$ 收敛，且其极限既等于 $f(x)$ 又等于 $g(x)$，故 $f(x)=g(x)$。因此 $\{x\in X\mid f(x)\neq g(x)\}\subseteq N\in\mathcal{N}_\mu$，即 $f=g$ $\mu$--a.e.。
\end{proof}

为后续参考，我们证明以下技术引理。

% SUPPLEMENT: parent=lemma:ch01-穷竭; kind=motivation; 该对象的动机与条件作用见紧随其后的注记
\begin{lemma}[耗竭引理]\label{lemma:ch01-穷尽引理}\label{lemma:ch01-穷竭}
设 $(X,\mathcal{A},\mu)$ 为一有限测度空间。设 $P=P(x)$（$x\in X$）为一性质，满足：
\begin{enumerate}[(i)]
\item $P$ 在 $X$ 的每个 $\mu$--零测子集上成立；
\item 对任意 $E\in\mathcal{A}$ 且 $\mu(E)>0$，存在 $F\in\mathcal{A}$、$F\subseteq E$、$\mu(F)>0$，使得 $P$ 在 $F$ 上成立。
\end{enumerate}
则存在不交序列 $\{E_k\}\subseteq\mathcal{A}$ 使得 $\bigcup_{k=1}^\infty E_k=X$ 且 $P$ 在每个 $E_k$ 上成立。
\end{lemma}

% SUPPLEMENT: parent=lemma:ch01-穷尽引理; kind=proof; 原证明断言 $\bigcup_{n=1}^kF_n$ 上 $P$ 成立，这需要 $\{E\mid P\text{ 在 }E\text{ 上成立}\}$ 对有限并封闭，而该封闭性并非假设 (i)(ii) 的推论；此处保留上确界论证的思路，改用逐次取上确界并拼接互不相交的集合，只用到 (i)(ii)
\begin{proof}
令 $\Sigma\triangleq\{E\in\mathcal{A}\mid P\text{ 在 }E\text{ 上成立}\}$，$c\triangleq\sup_{E\in\Sigma}\mu(E)$。由 $\mu(X)<\infty$ 得 $c\leqslant\mu(X)<\infty$。

现归纳地构造 $\Sigma$ 中两两不交的集合。令 $L_0\triangleq X$。若 $L_{k-1}$ 已定义，令
\begin{align*}
c_k\triangleq\sup\{\mu(E)\mid E\in\Sigma,\ E\subseteq L_{k-1}\},
\end{align*}
并区分两种情形。若 $c_k>0$，取 $F_k\in\Sigma$、$F_k\subseteq L_{k-1}$，使得 $\mu(F_k)\geqslant c_k/2$（由 $c_k$ 为上确界且 $c_k>0$ 可取到），并令 $L_k\triangleq L_{k-1}\setminus F_k$。若 $c_k=0$，则构造终止。

先说明若构造在 $k$ 步终止（即 $c_k=0$），则 $\mu(L_{k-1})=0$：否则由 (ii) 存在 $F\in\mathcal{A}$、$F\subseteq L_{k-1}$、$\mu(F)>0$，使 $P$ 在 $F$ 上成立，即 $F\in\Sigma$，于是 $c_k\geqslant\mu(F)>0$，与 $c_k=0$ 矛盾。

其次说明若构造不终止，则 $c_k\to0$。由 $F_j\subseteq L_{j-1}$ 知 $\{F_j\}$ 两两不交，故 $\sum_{j=1}^\infty\mu(F_j)\leqslant\mu(X)<\infty$；又当 $c_j>0$ 时 $\mu(F_j)\geqslant c_j/2$，即 $c_j\leqslant2\mu(F_j)$，于是 $\sum_{j=1}^\infty c_j\leqslant2\sum_{j=1}^\infty\mu(F_j)<\infty$，特别地 $c_k\to0$。

现设构造不终止，并令 $L\triangleq\bigcap_{k=0}^\infty L_k=X\setminus\bigcup_{k=1}^\infty F_k$。若 $\mu(L)>0$，则由 (ii) 存在 $G\in\mathcal{A}$、$G\subseteq L$、$\mu(G)>0$，使 $P$ 在 $G$ 上成立。则 $G\in\Sigma$ 且 $G\subseteq L_{k-1}$ 对每个 $k$ 成立，故 $\mu(G)\leqslant c_k$ 对每个 $k$ 成立，于是 $\mu(G)\leqslant\lim_{k\to\infty}c_k=0$，矛盾。故 $\mu(L)=0$，从而由 (i) 知 $P$ 在 $L$ 上成立。

于是 $\{F_1,F_2,\cdots\}$ 与 $L$ 是 $\mathcal{A}$ 中两两不交的集合，$P$ 在每一个上成立，且它们的并为 $X$。若构造在 $k$ 步终止，则 $\{F_1,\cdots ,F_{k-1}\}$ 与 $L_{k-1}$ 具有同样的性质（此时 $\mu(L_{k-1})=0$）。把该有限或可数族按任意方式枚举为 $\{E_k\}$（不足时补 $\varnothing$，$P$ 在 $\varnothing$ 上成立），即得所需的不交序列。
\end{proof}

% SUPPLEMENT: parent=lemma:ch01-穷尽引理; kind=motivation; 说明引理的作用（把全局性质化归为可数块上的局部性质）、两条件的分工、有限性假设的必要性，以及条件 (ii) 与"$P$ 对有限并封闭"的关系
\begin{remark}\phantomsection\label{remark:ch01-穷尽引理注}
\begin{enumerate}[(i)]
\item 引入动机。这是测度论中的标准"耗竭"论证：若每个正测集都含有一个正测的"好"子集，则整个空间可分成可数多个"好"块。它把关于全局的断言化归为关于可数多块的分段断言，因而是向量测度的 Radon--Nikodým 定理与各类分解定理的标准工具。
\item 两条件的分工。(i) 保证零测的"剩余部分"可以并入分割而不破坏性质；(ii) 是唯一用到 $\mu(E)>0$ 的地方，它保证只要还有正测度的剩余，就能取出新的"好"块，从而上确界不能停在 $\mu(X)$ 之下。
\item 假设的作用。$\mu(X)<\infty$ 在证明中两处使用：一是 $c\leqslant\mu(X)<\infty$ 保证上确界有限，从而可按 $\mu(F_k)\geqslant c_k/2$ 选出 $F_k$；二是 $\sum_j\mu(F_j)\leqslant\mu(X)<\infty$ 保证 $\sum_jc_j<\infty$，进而 $c_k\to0$，这是余集 $L$ 为零测集的关键。缺少有限性时该上确界论证失效，故引理只对有限测度空间陈述；对一般的 $\sigma$--有限测度，应先作 $\sigma$--有限的分解再逐块应用。
\item 边界情形。条件 (ii) 只要求 $P$ 在某个正测子集上成立，不要求 $P$ 对有限并封闭（这是与本条证明所要回避的困难相应的一般情形）；若 $P$ 对有限并封闭，则可直接取 $\bigcup_{n=1}^kF_n$ 并按更短的论证完成。若 $X$ 本身零测，则由 (i) 取 $E_1=X$ 即可。
\end{enumerate}
\end{remark}

最后我们证明关于 Radon 测度支集的如下刻画。

\begin{lemma}\label{lemma:ch01-支集刻画}
设 $(X,\mathcal{T})$ 为一 Hausdorff 空间，$\mu\in\mathcal{R}^+(X)$，$\mathcal{K}$ 为紧子集族。设 $A\in\mathcal{T}$ 为包含一切开 $\mu$--零测集的开 $\mu$--零测集（其存在性在证明中给出）。则 $\operatorname{supp}\mu=A^c$，即支集恰为一切开 $\mu$--零测集之并的补集。
\end{lemma}

% SUPPLEMENT: parent=lemma:ch01-支集刻画; kind=proof; 原证明中"$\mu(K)=0$"与"由内正则性得 $\mu(A)=0$"两处未展开，此处补出紧集被有限个零测开集覆盖的细节
\begin{proof}
先证最大的开 $\mu$--零测集存在。令 $\mathcal{N}_\mu\cap\mathcal{T}\equiv\{A_j\}_{j\in J}$，$A\triangleq\bigcup_{j\in J}A_j$。则 $A\in\mathcal{T}$ 且 $A\supseteq A_j$ 对一切 $j\in J$ 成立。若 $\mathcal{N}_\mu\cap\mathcal{T}=\{\varnothing \}$，则 $A=\varnothing$ 且结论成立（$\mu(A)=0$）。否则设 $K\subseteq A$ 为紧集。由于 $\{A_j\}_{j\in J}$ 是 $K$ 的开覆盖，存在 $A_{j_1},\cdots ,A_{j_n}$ 使 $K\subseteq\bigcup_{i=1}^nA_{j_i}$；由\cref{proposition:ch01-测度性质}(iii) 与 $\mu(A_{j_i})=0$ 得 $\mu(K)\leqslant\sum_{i=1}^n\mu(A_{j_i})=0$。由 $K$ 的任意性与 $\mu$ 的内正则性，
\begin{align*}
\mu(A)=\sup\{\mu(K)\mid K\subseteq A,\ K\in\mathcal{K}\}=0,
\end{align*}
故 $A$ 是包含一切开 $\mu$--零测集的开 $\mu$--零测集。

再证 $\operatorname{supp}\mu=A^c$。设 $x\in[\operatorname{supp}\mu]^c$。则存在 $x$ 的开邻域 $\widetilde{A}$ 使 $\mu(\widetilde{A})=0$（否则 $x$ 的每个邻域测度为正，与 $x\notin\operatorname{supp}\mu$ 矛盾）；特别地 $\widetilde{A}\in\mathcal{N}_\mu\cap\mathcal{T}$，故 $x\in\widetilde{A}\subseteq A$。反之，设 $x\in A$。则 $A$ 是 $x$ 的 $\mu$--零测开邻域，特别地 $\mu(A)=0$，故 $x$ 不满足\cref{definition:ch01-Borel测度}(g) 中"每个邻域测度为正"的条件，即 $x\notin\operatorname{supp}\mu$。故 $[\operatorname{supp}\mu]^c=A$，即 $\operatorname{supp}\mu=A^c$。
\end{proof}

% SUPPLEMENT: parent=lemma:ch01-支集刻画; kind=motivation; 说明支集与最大开零测集的关系、内正则性的作用、以及缺少 Radon 假设时结论失效的例子
\begin{remark}\phantomsection\label{remark:ch01-支集刻画注}
\begin{enumerate}[(i)]
\item 引入动机。\cref{definition:ch01-Borel测度}(g) 用"每个邻域测度为正"定义支集，引理则把它刻画为最大开零测集的补集：支集是使 $\mu$ 集中的最小闭集，测度在支集之外"完全消失"。这一刻画在讨论测度的支集、narrow 收敛与集中性质时是基本的。
\item 假设的作用。$\mu$ 的内正则性只在证明的第一段使用：需由 $\mu(A)=\sup\{\mu(K)\mid K\subseteq A,\ K\in\mathcal{K}\}$ 与"每个紧集 $K\subseteq A$ 都被有限个零测开集覆盖"推出 $\mu(A)=0$。局部有限性保证 $\mu$ 在紧集上有限，从而使该有限求和有意义。Hausdorff 假设属于 Radon 测度的论域设定（见\cref{definition:ch01-Borel测度}），在本证明中不额外使用。
\item 边界情形。缺少内正则性时，"最大开零测集"甚至可以不存在。取 $X=\mathbb{R}$ 与计数测度 $\mu$（每个单点集既开又零测），则每个 $x$ 都有零测开邻域 $\{x\}$，故按\cref{definition:ch01-Borel测度}(g) 中的定义式，支集为 $\mathbb{R}$；但一切开零测集之并为 $\mathbb{R}$，它并不零测，故不存在最大的开零测集。此例中 $\mu$ 不局部有限（每个非空开集测度为 $\infty$），故不满足引理的 Radon 假设。又若 $\mu\equiv0$，则 $A=X$、$\operatorname{supp}\mu=\varnothing$，与"支集可以为空"一致。
\end{enumerate}
\end{remark}

\end{document}
